%!TEX root = ../../../thesis.tex


\IN{Boolean transition matrix}
\label{def:transitionmtx} 

Let $\semiautomat{A}{Q}{\Sigma}{\delta}$ be an $n$-state automaton.
As shown in \cref{prelim:subsec:transsemigrp}, letters and words induce transformations
on the states of $\mathcal A$, which can be represented
by transition matrices similar as in \cref{prelim:eqn:trans_mtx}.
More precisely, we have to fix some ordering on the states $Q = \{q_1, \ldots, q_n\}$.
Then, to every letter $a \in \Sigma$ we associate a 
\emph{(Boolean) transition matrix}
$[a]_{\mathcal A} = (m^{(a)}_{ij})_{i,j \in \{1,\ldots,n\}}$ by setting
\[
m^{(a)}_{ij} = \begin{cases}
	1 & \text{if } \applytrans{\delta}{q_i}{a} = q_j; \\
	0 & \text{if } \applytrans{\delta}{q_i}{a} \ne q_j.
\end{cases} 
\]
%So, via the ordering $\{q_1, \ldots, q_n\}$
%we identify each state uniquely with a number in $\{1,\ldots,n\}$,
%which correspond to the rows and columns of the transition matrices.
Intuitively, state $q_i$ is mapped to state $q_j$ by $a \in \Sigma$ \todo{das vorne bei transitions so schon sagen?}
if and only if in the $i$-th row of~$[a]_{\mathcal A}$
the unique one appears at position $j$, i.e., in the $j$-th column of $[a]_{\mathcal A}$.
Inductively\footnote{Alternatively, we can apply the Universal Mapping Property (Theorem~\ref{prelim:prop:UMP})
	to the monoid of all $n \times n$ matrices and the mapping $[\cdot]_{\mathcal A}$ from letters into this monoid
	to obtain the unique morphism from $\Sigma^*$ into the set of $n \times n$ Boolean matrices.}, 
we associate with every word a transition matrix
in the following way: we let $[\varepsilon]_{\mathcal A}$ be the identity matrix $I_n$
and, for $u \in \Sigma^*$ and $a \in \Sigma$, set $[ua]_{\mathcal A} = [u]_{\mathcal A}[a]_{\mathcal A}$, i.e., the matrix product 
of~$[u]_{\mathcal A}$ and~$[a]_{\mathcal A}$.
\Cref{prelim:fig:transitionmtxCerny} gives transition matrices for the \v{C}ern\'y automata $\mathcal C_n$
from \cref{intro:sync_aut:cerny}.

\begin{figure}[ht]
	\centering
	\[
	[b]_{\mathcal C_n} =
	\left( {\begin{array}{ccccc}
			0 & 1 &  0 & \cdots & 0\\
			0 & 0 &  1 & \cdots & 0\\
			\vdots & \vdots & \vdots & \ddots & \vdots\\
			0 & 0 & 0 & \cdots & 1 \\
			1 & 0 & 0 & \cdots & 0
	\end{array} } \right)
	\qquad 
	[a]_{\mathcal C_n} =
	\left( {\begin{array}{cccccc}
			0 & 1 &  0 & \cdots & 0 & 0 \\
			0 & 1 &  0 & \cdots & 0 & 0 \\
			\vdots & \vdots & \vdots & \ddots & \vdots\\
			0 & 0 & 0 & \cdots & 1 & 0 \\
			0 & 0 & 0 & \cdots & 0 & 1
	\end{array} } \right)
	\]
	\caption[The transition matrices for the \v{C}ern\'y automata $\mathcal C_n$]{
		The transition matrices for the \v{C}ern\'y automata $\mathcal C_n$
		from \SEC~\ref{intro:sync_aut:cerny} (\FIG~\ref{intro:fig:cerny})
		with the state identification $q_1 = 0, q_2 = 1, \ldots, q_n = n - 1$.
	}
	\label{prelim:fig:transitionmtxCerny}
\end{figure} 

When there is no ambiguity to which automaton we refer,
we write $[u]$ instead of $[u]_{\mathcal A}$ for $u \in \Sigma^*$.
Also, when the specific ordering of the states is not important, we
assume some order to be taken without explicitly mentioning it.
% do not mention it and assume some order to be taken.
